\label{sec:methodology}

\subsection{Algorithmic District Generation}

We use the \texttt{redist} recursive-bisection system~\citep{dellaLibera2026recursive}
to generate congressional maps for all 50 states using 2020 census tract data.
The algorithm partitions census tracts into $k$ districts---where $k$ is the
state's congressional apportionment---by minimizing edge-cut in the tract
adjacency graph subject to population balance within $\pm 0.5\%$ of the
equal-population target.  The algorithm receives only population counts and
geographic adjacency data; it has no access to election results, voter
registration records, or any partisan information.

For states with a single at-large district ($k=1$: Alaska, Delaware, Montana,
North Dakota, South Dakota, Vermont, Wyoming), we treat the entire state as a
single district.  For multi-district states, we run the full recursive
bisection pipeline as described in~\citet{dellaLibera2026recursive}.

\subsection{Enacted Districts}

Enacted district boundaries for the 117th Congress (2021 cycle, based on 2020
census) are taken from the U.S. Census Bureau's TIGER/Line shapefiles.  These
represent the maps actually used for the 2022 midterm elections.

\subsection{Vote-Share Reaggregation}

We reaggregate 2020 presidential precinct-level results to both algorithmic
and enacted districts using area-weighted interpolation.  Where precinct
boundaries do not align exactly with census tract boundaries, we apportion
votes proportionally to the overlapping area.  This standard procedure
introduces small measurement error, but since both algorithmic and enacted
maps are subject to the same reaggregation process, the comparison is unbiased.

\subsection{Comparison Framework}

For each state, we compute:
\begin{enumerate}
  \item The number of competitive districts ($|D_i - 0.5| < 0.10$) under the
    algorithmic map.
  \item The number of competitive districts under the enacted map.
  \item The difference (algorithmic minus enacted).
  \item The partisan split of competitive districts under each map.
\end{enumerate}

National totals are aggregated across all 50 states.  We report the
distribution of competitive district counts and the aggregate partisan split.

\subsection{Longitudinal Projection}

To assess whether competitive algorithmic districts remain competitive as vote
shares drift, we apply uniform partisan shifts of $\pm 1$, $\pm 2$, $\pm 3$,
and $\pm 4$ percentage points to the 2020 presidential baseline.  A district
that remains within the competitive window under the full range of $\pm 4$
point shifts is classified as \emph{durably competitive}.
